Large sparse linear systems arise throughout scientific computing, engineering simulation, and numerical modelling. These systems are often written in the form
\begin{equation}
    Ax
    =
    b,
    \label{eq:matrix-equation}
\end{equation}
where \(A\) is a large sparse matrix, \(b\) is a known right-hand side, and \(x\) is the unknown solution vector. For problems of realistic size, direct methods such as matrix factorisation can require too much memory or computation, so iterative methods are commonly used instead.

GMRES, the Generalized Minimal Residual method, is one of the standard Krylov subspace methods for solving non-symmetric linear systems \cite{saad-2003,trefethen-1997}. Rather than solving \Cref{eq:matrix-equation} directly, GMRES builds a sequence of approximate solutions from Krylov subspaces generated by the matrix \(A\) and the initial residual. At each iteration, it chooses the approximation whose residual is minimal over the current subspace. This residual-minimising property makes GMRES a robust and widely used method for difficult sparse systems.

The difficulty addressed in this thesis comes from running GMRES on modern high-performance computing systems. In a parallel implementation, the matrix and vectors are distributed across multiple processes. Each process can perform much of the arithmetic on its own local data, but operations such as inner products and vector norms require information from every process. For example, after the current Arnoldi vector \(w=Aq_i\) has been formed, GMRES needs the projection coefficient \(h_{j,i}=q_j^T w\). If process \(p\) owns the index set \(I_p\), it can only compute its local contribution
\begin{equation}
    \alpha_p
    =
    \sum_{k\in I_p} q_{j,k}w_k.
    \label{eq:local-dot-product}
\end{equation}
The global coefficient is obtained by summing these local contributions,
\begin{equation}
    h_{j,i}
    =
    q_j^T w
    =
    \sum_p \alpha_p.
    \label{eq:global-dot-product}
\end{equation}
This sum is usually performed using a collective operation such as \texttt{MPI\_Allreduce}. It is necessary for correctness, because no single process has the whole vector, but it introduces communication and synchronization before the algorithm can continue.

\begin{figure}[h]
    \centering
    \resizebox{0.95\textwidth}{!}{%
    \begin{tikzpicture}[
        font=\small,
        >=Stealth,
        process/.style={draw=TCDGrey75, fill=TCDGrey10, minimum width=0.95cm, minimum height=0.55cm, align=center},
        vector/.style={draw=TCDBlue, fill=TCDBlue10, minimum width=1.35cm, minimum height=0.55cm, align=center},
        work/.style={draw=TCDGrey, fill=white, minimum width=1.45cm, minimum height=0.55cm, align=center},
        comm/.style={draw=TCDRed, fill=TCDRed10, thick, minimum width=2.65cm, minimum height=1.35cm, align=center},
        result/.style={draw=TCDGreen, fill=TCDGreen10, thick, minimum width=3.0cm, minimum height=0.8cm, align=center},
        update/.style={draw=TCDGrey, fill=TCDGrey10, minimum width=3.35cm, minimum height=0.95cm, align=center},
        heading/.style={font=\bfseries, align=center},
        arrow/.style={->, thick}
    ]
        \node[font=\bfseries, anchor=west] at (0.1,4.95) {One Arnoldi dot product in distributed GMRES};

        \node[heading] at (2.55,4.35) {local vector slices};
        \node[heading] at (5.8,4.35) {partial sums};
        \node[heading] at (8.8,4.35) {global reduction};
        \node[heading] at (12.15,4.35) {shared scalar};

        \node[comm] (reduce) at (8.8,2.45) {\texttt{MPI\_Allreduce}\\[2pt]sum and share\\[2pt]\(h_{j,i}=\sum_p\alpha_p\)};
        \node[result] (coef) at (12.15,3.05) {\(h_{j,i}=q_j^T w\)};
        \node[update] (upd) at (12.15,1.55) {all processes use \(h_{j,i}\)\\to update their local slice};

        \foreach \p/\y in {0/3.6,1/2.85,2/2.1,3/1.35} {
            \node[process] at (0.45,\y) {proc. \(\p\)};
            \node[vector] (q\p) at (2.0,\y) {\(q_j^{(\p)}\)};
            \node[vector] (w\p) at (3.45,\y) {\(w^{(\p)}\)};
            \node[work] (a\p) at (5.8,\y) {\(\alpha_{\p}\)};

            \draw[arrow] (w\p.east) -- (a\p.west);
            \draw[arrow] (a\p.east) -- (reduce.west);
        }

        \draw[arrow] (reduce.east) -- (coef.west);
        \draw[arrow] (coef.south) -- (upd.north);
    \end{tikzpicture}}
    \caption{In distributed Arnoldi, each process stores only slices of the Krylov basis vector \(q_j\) and the current vector \(w\). It can form only a local partial sum \(\alpha_p\). A global reduction sums these partial values and returns the shared coefficient \(h_{j,i}\), which every process then uses in its local orthogonalisation update.}
    \label{fig:gmres-communication}
\end{figure}

Communication-avoiding Krylov methods attempt to reduce this bottleneck by reorganising the algorithm so that fewer synchronization operations are required. In particular, \(s\)-step GMRES variants generate several Krylov vectors at a time,
\begin{equation}
    K_s
    =
    \left[
        q_i,\,
        Aq_i,\,
        \ldots,\,
        A^{s-1}q_i
    \right],
    \label{eq:s-step-krylov-block}
\end{equation}
increasing the amount of local computation performed between communication phases. However, forming this block is only part of the problem. The vectors in \(K_s\) must still be orthogonalised against the existing Krylov basis and against each other. This requires block inner products such as \(Q_i^T K_s\) and \(K_s^T K_s\), whose entries are also global sums over all processes. If those products are computed one column at a time, the method simply reintroduces the many reductions that \(s\)-step GMRES was meant to avoid. The block orthogonalisation therefore also has to be communication avoiding: the aim is to replace many small scalar reductions with fewer collective operations on small dense matrices, while still controlling the loss of orthogonality caused by finite precision arithmetic.

\begin{figure}[h]
    \centering
    \resizebox{0.95\textwidth}{!}{%
    \begin{tikzpicture}[
        font=\small,
        >=Stealth,
        local/.style={draw=TCDBlue, fill=TCDBlue10, minimum width=1.35cm, minimum height=0.7cm, align=center},
        block/.style={draw=TCDBlue, fill=TCDBlue10, thick, minimum width=4.2cm, minimum height=1.35cm, align=center},
        comm/.style={draw=TCDRed, fill=TCDRed10, thick, minimum width=1.55cm, minimum height=0.85cm, align=center},
        ca/.style={draw=TCDRed, fill=TCDRed10, thick, minimum width=3.15cm, minimum height=1.3cm, align=center},
        result/.style={draw=TCDGreen, fill=TCDGreen10, thick, minimum width=2.55cm, minimum height=0.9cm, align=center},
        note/.style={draw=TCDGrey75, fill=TCDGrey10, minimum width=3.0cm, minimum height=0.7cm, align=center},
        heading/.style={font=\bfseries, anchor=west},
        arrow/.style={->, thick}
    ]
        \node[heading] at (0.1,4.95) {Classical Arnoldi: many synchronization points};

        \node[local] (c1) at (0.95,4.05) {make\\\(Aq_i\)};
        \node[comm] (r1) at (2.85,4.05) {reduce};
        \node[local] (c2) at (4.75,4.05) {make\\\(Aq_{i+1}\)};
        \node[comm] (r2) at (6.65,4.05) {reduce};
        \node[local] (c3) at (8.55,4.05) {make\\\(Aq_{i+2}\)};
        \node[comm] (r3) at (10.45,4.05) {reduce};
        \node[note] (syncs) at (10.45,3.05) {each red box is a global wait};

        \draw[arrow] (c1.east) -- (r1.west);
        \draw[arrow] (r1.east) -- (c2.west);
        \draw[arrow] (c2.east) -- (r2.west);
        \draw[arrow] (r2.east) -- (c3.west);
        \draw[arrow] (c3.east) -- (r3.west);
        \draw[arrow] (r3.south) -- (syncs.north);

        \node[heading] at (0.1,2.35) {\(s\)-step GMRES: build a block, then orthogonalise the block};

        \node[block] (ks) at (2.2,1.0) {};
        \draw[TCDBlue] (1.0,0.325) -- (1.0,1.675);
        \draw[TCDBlue] (2.0,0.325) -- (2.0,1.675);
        \draw[TCDBlue] (2.8,0.325) -- (2.8,1.675);
        \node at (0.55,1.0) {\(q_i\)};
        \node at (1.5,1.0) {\(Aq_i\)};
        \node at (2.4,1.0) {\(\cdots\)};
        \node at (3.55,1.0) {\(A^{s-1}q_i\)};

        \node[ca] (orth) at (6.4,1.0) {CA block\\orthogonalisation\\[2pt]\(Q_i^T K_s,\ K_s^T K_s\)};
        \node[result] (qs) at (10.15,1.0) {orthonormal\\block \(Q_s\)};
        \node[note] (few) at (6.4,-0.45) {fewer reductions, but still global};

        \draw[arrow] (ks.east) -- (orth.west);
        \draw[arrow] (orth.east) -- (qs.west);
        \draw[arrow] (orth.south) -- (few.north);
    \end{tikzpicture}}
    \caption{The \(s\)-step idea groups several Krylov-vector generations into a block, reducing the number of times the algorithm has to stop for communication. The block still has to be orthogonalised using global inner products, so communication avoidance must also be built into the block orthogonalisation stage.}
    \label{fig:s-step-communication-avoidance}
\end{figure}

In principle, this can reduce communication latency by a factor related to the step size \(s\), but in finite precision the step size cannot be chosen arbitrarily large without risking numerical instability. The adaptive \(s\)-step GMRES algorithm proposed by Xu, Alonso, and Darve uses this idea while dynamically controlling the block size in order to maintain stability \cite{xu-2023}.

The aim of this thesis is to implement and study a communication-avoiding \(s\)-step GMRES solver in a parallel computing environment. A baseline classical GMRES implementation is used as a reference point, so that the communication-avoiding method can be compared against the standard algorithm in terms of convergence behaviour, numerical stability, and parallel performance. The work therefore sits between numerical linear algebra and high-performance computing: the mathematical question is whether the Krylov basis remains stable and accurate, while the computational question is whether reducing synchronization improves scalability.

The remainder of the thesis develops this study step by step. \Cref{sec:background} introduces the numerical linear algebra and parallel computing concepts needed for the project, including Krylov subspaces, GMRES, orthogonalisation, and communication in distributed-memory systems. \Cref{sec:implementation} then describes the solver implementations, including sparse matrix operations, MPI communication, and the testing strategy. \Cref{sec:results} presents convergence and performance results, comparing the classical and communication-avoiding approaches. Finally, \Cref{sec:conclusion} summarises the main findings and discusses possible directions for future work.
